
%% bare_lsens.tex
%% V1.0
%% 2017/07/12
%% by Michael Shell
%% see http://www.michaelshell.org/
%% for current contact information.
%%
%% This is a skeleton file demonstrating the use of IEEE_lsens.cls for
%% an IEEE Sensors Letters paper.
%%
%% Support sites:
%% -- LaTeX Related --
%% http://www.michaelshell.org/tex/ieeetran/
%% http://www.ctan.org/pkg/ieeetran/
%%
%% -- IEEE/Editor/Submission Related --
%% http://www.ieee-sensors.org/
%% http://www.ieee.org/


%%*************************************************************************
%% Legal Notice:
%% This code is offered as-is without any warranty either expressed or
%% implied; without even the implied warranty of MERCHANTABILITY or
%% FITNESS FOR A PARTICULAR PURPOSE! 
%% User assumes all risk.
%% In no event shall the IEEE or any contributor to this code be liable for
%% any damages or losses, including, but not limited to, incidental,
%% consequential, or any other damages, resulting from the use or misuse
%% of any information contained here.
%%
%% All comments are the opinions of their respective authors and are not
%% necessarily endorsed by the IEEE.
%%
%% This work is distributed under the LaTeX Project Public License (LPPL)
%% ( http://www.latex-project.org/ ) version 1.3, and may be freely used,
%% distributed and modified. A copy of the LPPL, version 1.3, is included
%% in the base LaTeX documentation of all distributions of LaTeX released
%% 2003/12/01 or later.
%% Retain all contribution notices and credits.
%% ** Modified files should be clearly indicated as such, including  **
%% ** renaming them and changing author support contact information. **
%%*************************************************************************


% *** Authors should verify (and, if needed, correct) their LaTeX system  ***
% *** with the testflow diagnostic prior to trusting their LaTeX platform ***
% *** with production work. The IEEE's font choices and paper sizes can   ***
% *** trigger bugs that do not appear when using other class files.       ***
% The testflow support page is at:
% http://www.michaelshell.org/tex/testflow/


% Typically, there isn't any need for class options with IEEE Sensors
% Letters papers. The default, and only, text size is 9pt.
\documentclass{IEEE_lsens}
%
% If IEEE_lsens.cls has not been installed into the LaTeX system files,
% manually specify the path to it like:
% \documentclass{../sty/IEEE_lsens}


% *** Do not adjust lengths that control margins, column widths, etc. ***
% *** Do not use packages that alter fonts (such as pslatex).         ***
% There should be no need to do such things with IEEE_lsens.
% (Unless specifically instructed to do so by the IEEE Sensors Letters
% editors, of course.)


% Some very useful LaTeX packages include:
% (uncomment the ones you want to load)


% *** MISC UTILITY PACKAGES ***
%
%\usepackage{ifpdf}
% Heiko Oberdiek's ifpdf.sty is very useful if you need conditional
% compilation based on whether the output is pdf or dvi.
% usage:
% \ifpdf
%   % pdf code
% \else
%   % dvi code
% \fi
% The latest version of ifpdf.sty can be obtained from:
% http://www.ctan.org/pkg/ifpdf
% Also, note that IEEEtran.cls V1.7 and later provides a builtin
% \ifCLASSINFOpdf conditional that works the same way.
% When switching from latex to pdflatex and vice-versa, the compiler may
% have to be run twice to clear warning/error messages.


% The textcomp package can be loaded to get better \textcopyright and
% \textregistered symbols (as is used in the footer of the title page
% of LSENS papers).
\usepackage{textcomp}
\usepackage{graphicx}
\usepackage{siunitx}






% *** CITATION PACKAGES ***
%
\usepackage[noadjust]{cite}
% cite.sty was written by Donald Arseneau
% V1.6 and later of IEEEtran pre-defines the format of the cite.sty package
% \cite{} output to follow that of the IEEE. Loading the cite package will
% result in citation numbers being automatically sorted and properly
% "compressed/ranged". e.g., [1], [9], [2], [7], [5], [6] without using
% cite.sty will become [1], [2], [5]--[7], [9] using cite.sty. cite.sty's
% \cite will automatically add leading space, if needed. Use cite.sty's
% noadjust option (cite.sty V3.8 and later) if you want to turn this off
% such as if a citation ever needs to be enclosed in parenthesis.
% cite.sty is already installed on most LaTeX systems. Be sure and use
% version 5.0 (2009-03-20) and later if using hyperref.sty.
% The latest version can be obtained at:
% http://www.ctan.org/pkg/cite
% The documentation is contained in the cite.sty file itself.






% *** GRAPHICS RELATED PACKAGES ***
%
\ifCLASSINFOpdf
  % \usepackage[pdftex]{graphicx}
  % declare the path(s) where your graphic files are
  % \graphicspath{{../pdf/}{../jpeg/}}
  % and their extensions so you won't have to specify these with
  % every instance of \includegraphics
  % \DeclareGraphicsExtensions{.pdf,.jpeg,.png}
\else
  % or other class option (dvipsone, dvipdf, if not using dvips). graphicx
  % will default to the driver specified in the system graphics.cfg if no
  % driver is specified.
  % \usepackage[dvips]{graphicx}
  % declare the path(s) where your graphic files are
  % \graphicspath{{../eps/}}
  % and their extensions so you won't have to specify these with
  % every instance of \includegraphics
  % \DeclareGraphicsExtensions{.eps}
\fi
% graphicx was written by David Carlisle and Sebastian Rahtz. It is
% required if you want graphics, photos, etc. graphicx.sty is already
% installed on most LaTeX systems. The latest version and documentation
% can be obtained at: 
% http://www.ctan.org/pkg/graphicx
% Another good source of documentation is "Using Imported Graphics in
% LaTeX2e" by Keith Reckdahl which can be found at:
% http://www.ctan.org/pkg/epslatex
%
% latex, and pdflatex in dvi mode, support graphics in encapsulated
% postscript (.eps) format. pdflatex in pdf mode supports graphics
% in .pdf, .jpeg, .png and .mps (metapost) formats. Users should ensure
% that all non-photo figures use a vector format (.eps, .pdf, .mps) and
% not a bitmapped formats (.jpeg, .png). The IEEE frowns on bitmapped formats
% for line art (i.e., plots, graphs, etc.) as they result in 
% "jaggedy"/blurry rendering of lines and letters as well as large
% increases in file sizes.
%
% You can find documentation about the pdfTeX application at:
% http://www.tug.org/applications/pdftex





% *** MATH PACKAGES ***

\usepackage[T1]{fontenc} % optional enhanced font encoding
\usepackage{amsmath}
% A popular package from the American Mathematical Society that provides
% many useful and powerful commands for dealing with mathematics.
%
% Note that the amsmath package sets \interdisplaylinepenalty to 10000
% thus preventing page breaks from occurring within multiline equations. Use:
\interdisplaylinepenalty=2500
% after loading amsmath to restore such page breaks as IEEEtran.cls normally
% does. amsmath.sty is already installed on most LaTeX systems. The latest
% version and documentation can be obtained at:
% http://www.ctan.org/pkg/amsmath

\usepackage[cmintegrals]{newtxmath}
% A Times compatible math font is required for IEEE Sensors Letters.
% Michael Sharpe's freely available newtxmath package (version 1.451,
% July 28, 2015 or later) is recommended. The cmintegrals option, which
% IEEE_lsens sets as a default upon loading newtxmath, is/was needed to
% obtain the specific style of integral symbol used by the IEEE Sensors
% Letters. However, as of version 1.5 of the newtxmath package, released
% on 2016/08/12, the correct integral symbol is now invoked by default
% and the cmintegrals option is no longer needed and is silently ignored.
% The latest version and documentation can be obtained at:
% http://www.ctan.org/pkg/newtx
\usepackage{bm}
% support for selective bold math 
% The latest version and documentation can be obtained at:
% http://www.ctan.org/pkg/bm




% *** SPECIALIZED LIST PACKAGES ***
%
%\usepackage{algorithmic}
% algorithmic.sty was written by Peter Williams and Rogerio Brito.
% This package provides an algorithmic environment fo describing algorithms.
% You can use the algorithmic environment in-text or within a figure
% environment to provide for a floating algorithm. Do NOT use the algorithm
% floating environment provided by algorithm.sty (by the same authors) or
% algorithm2e.sty (by Christophe Fiorio) as the IEEE does not use dedicated
% algorithm float types and packages that provide these will not provide
% correct IEEE style captions. The latest version and documentation of
% algorithmic.sty can be obtained at:
% http://www.ctan.org/pkg/algorithms
% Also of interest may be the (relatively newer and more customizable)
% algorithmicx.sty package by Szasz Janos:
% http://www.ctan.org/pkg/algorithmicx




% *** ALIGNMENT PACKAGES ***
%
\usepackage{array}
% Frank Mittelbach's and David Carlisle's array.sty patches and improves
% the standard LaTeX2e array and tabular environments to provide better
% appearance and additional user controls. As the default LaTeX2e table
% generation code is lacking to the point of almost being broken with
% respect to the quality of the end results, all users are strongly
% advised to use an enhanced (at the very least that provided by array.sty)
% set of table tools. array.sty is already installed on most systems. The
% latest version and documentation can be obtained at:
% http://www.ctan.org/pkg/array


% IEEEtran contains the IEEEeqnarray family of commands that can be used to
% generate multiline equations as well as matrices, tables, etc., of high
% quality.




% *** SUBFIGURE PACKAGES ***
%\usepackage[caption=false,font=footnotesize,labelfont=sf,textfont=sf]{subfig}
% subfig.sty, written by Steven Douglas Cochran, is the modern replacement
% for subfigure.sty, the latter of which is no longer maintained and is
% incompatible with some LaTeX packages including fixltx2e. However,
% subfig.sty requires and automatically loads Axel Sommerfeldt's caption.sty
% which will override IEEEtran.cls' handling of captions and this will result
% in non-IEEE style figure/table captions. To prevent this problem, be sure
% and invoke subfig.sty's "caption=false" package option (available since
% subfig.sty version 1.3, 2005/06/28) as this is will preserve IEEE_lsens.cls
% handling of captions.
%
% The latest version and documentation of subfig.sty can be obtained at:
% http://www.ctan.org/pkg/subfig




% *** PDF, URL AND HYPERLINK PACKAGES ***
%
\usepackage{url}
% url.sty was written by Donald Arseneau. It provides better support for
% handling and breaking URLs. url.sty is already installed on most LaTeX
% systems. The latest version and documentation can be obtained at:
% http://www.ctan.org/pkg/url
% Basic use: \url{my_url_here}.
%
%
% NOTE: PDF hyperlink and bookmark features are not required in IEEE
%       papers and their use requires extra complexity and work.
\newcommand\MYhyperrefoptions{hypertexnames=true, bookmarks=true,
bookmarksnumbered=true, pdfpagemode={UseOutlines}, plainpages=false,
pdfpagelabels=true, colorlinks=true, linkcolor={black},
citecolor={black}, urlcolor={black}}
\ifCLASSINFOpdf
% \usepackage[\MYhyperrefoptions,breaklinks=true,pdftex]{hyperref}
\else
%  \usepackage[\MYhyperrefoptions,dvips]{hyperref}
%  \usepackage{breakurl}
\fi
%
% We will provide for these commands even if hyperref is not loaded to
% allow hyperref to be unloaded without have to delete any apprearances
% of these commands in the document.
\providecommand{\hypersetup}[1]{\relax}
\providecommand{\texorpdfstring}[2]{#1}
%
% We use \hypersetup instead of the package options for the
% PDF strings because there is a problem with using underscores
% with the package option approach. Also, the info that needs
% to be changed is all here. 
% *** IF USING HYPERREF BE SURE AND CHANGE THE EXAMPLE PDF ***
% *** TITLE/SUBJECT/AUTHOR/KEYWORDS INFO BELOW!!
%
\hypersetup{pdftitle={Bare Demo of IEEE\_lsens.cls for IEEE Sensors Letters},%<!CHANGE!
pdfsubject={Typesetting},%<!CHANGE!
pdfauthor={Michael D. Shell},%<!CHANGE!
pdfkeywords={Class, IEEE, IEEE\_lsens, IEEE Sensors Letters, LaTeX, Typesetting, TeX}}%<^!CHANGE!}
%
% One significant drawback of using hyperref under DVI output is that the
% LaTeX compiler cannot break URLs across lines or pages as can be done
% under pdfLaTeX's PDF output via the hyperref pdftex driver. This is
% probably the single most important capability distinction between the
% DVI and PDF output. Perhaps surprisingly, all the other PDF features
% (PDF bookmarks, thumbnails, etc.) can be preserved in
% .tex->.dvi->.ps->.pdf workflow if the respective packages/scripts are
% loaded/invoked with the correct driver options (dvips, etc.). 
% As most IEEE papers use URLs sparingly (mainly in the references), this
% may not be as big an issue as with other publications.
%
% That said, Vilar Camara Neto created his breakurl.sty package which
% permits hyperref to easily break URLs even in dvi mode.
% Note that breakurl, unlike most other packages, must be loaded
% AFTER hyperref. The latest version of breakurl and its documentation can
% be obtained at:
% http://www.ctan.org/pkg/breakurl
% breakurl.sty is not for use under pdflatex pdf mode.
%
% The advanced features offer by hyperref.sty are not required for IEEE
% submission, so users should weigh these features against the added
% complexity of use.
% The package options above demonstrate how to enable PDF bookmarks
% (a type of table of contents viewable in Acrobat Reader) as well as
% PDF document information (title, subject, author and keywords) that is
% viewable in Acrobat reader's Document_Properties menu. PDF document
% information is also used extensively to automate the cataloging of PDF
% documents. The above set of options ensures that hyperlinks will not be
% colored in the text and thus will not be visible in the printed page,
% but will be active on "mouse over". USING COLORS OR OTHER HIGHLIGHTING
% OF HYPERLINKS CAN RESULT IN DOCUMENT REJECTION BY THE IEEE, especially if
% these appear on the "printed" page. IF IN DOUBT, ASK THE RELEVANT
% SUBMISSION EDITOR. You may need to change/add the option for
% hypertexnames=false if you used duplicate equation numbers, etc., but
% this should not be needed in normal IEEE work.
% The latest version of hyperref and its documentation can be obtained at:
% http://www.ctan.org/pkg/hyperref



% correct bad hyphenation here
\hyphenation{op-tical net-works semi-conduc-tor}


\begin{document}


% The paper headers
% \markboth{Vol.~1, No.~3, May~2023}{0000000}


% article subject line 
\IEEELSENSarticlesubject{Sensor Applications}


% paper title
% Titles are generally capitalized except for words such as a, an, and, as,
% at, but, by, for, in, nor, of, on, or, the, to and up, which are usually
% not capitalized unless they are the first or last word of the title.
% Linebreaks \\ can be used within to get better formatting as desired.
% Do not put math or special symbols in the title.
%
\title{AlGaN/GaN Hall-Effect Sensor for In-Situ Magnetic Field Monitoring of the HSX Stellarator}


% author names and IEEE memberships
% note positions of commas and nonbreaking spaces ( ~ ) LaTeX will not break
% a structure at a ~ so this keeps an author's name from being broken across
% two lines.
% use \thanks{} to gain access to the first footnote area
% a separate \thanks must be used for each paragraph as LaTeX2e's \thanks
% was not built to handle multiple paragraphs
%
\author{\IEEEauthorblockN{Yiming~Zhao\IEEEauthorrefmark{1}, Wayne~Goodman\IEEEauthorrefmark{2}\IEEEauthorieeemembermark{1}, Thomas~Gallenberger\IEEEauthorrefmark{2}, Jasmine~M.~Cox\IEEEauthorrefmark{1}, Benedikt~Geiger\IEEEauthorrefmark{2}, and~Debbie~G.~Senesky\IEEEauthorrefmark{1,3}\IEEEauthorieeemembermark{2}}% <-this % stops a space
\IEEEauthorblockA{\IEEEauthorrefmark{1}Department of Electrical Engineering, Stanford University, Stanford, CA 94305 USA\\
\IEEEauthorrefmark{2}Department of Nuclear Engineering and Engineering Physics, University of Wisconsin-Madison, Madison, WI 53706 USA\\
\IEEEauthorrefmark{3}Department of Aeronautics and Astronautics, Stanford University, Stanford, CA 94305 USA\\
\IEEEauthorieeemembermark{1}Member, IEEE\\
\IEEEauthorieeemembermark{2}Senior Member, IEEE}%
% LSENS authors should provide a real e-mail address here.
\thanks{Corresponding authors: Yiming Zhao (e-mail: timzhao@stanford.edu), Debbie G. Senesky (e-mail: dsenesky@stanford.edu).}% <-this % stops a space
\thanks{Associate Editor: }%
\thanks{Digital Object Identifier }}
%
% note the % following lines that end in } - 
% these prevent an unwanted space from occurring between the objects.
% i.e., if you had this:
% 
% \author{....lastname \thanks{...} \thanks{...} }
%                     ^------------^------------^----Do not want these spaces!
%
% a space would be appended to the last name and could cause every name on that
% line to be shifted left slightly. This is one of those "LaTeX things". For
% instance, "\textbf{A} \textbf{B}" will typeset as "A B" not "AB". To get
% "AB" then you have to do: "\textbf{A}\textbf{B}"
% \thanks is no different in this regard, so shield the last } of each \thanks
% that ends a line with a % and do not let a space in before the next \thanks.
% For what it is worth, this is a minor point as most people would not even
% notice if the said evil space somehow managed to creep in.


% Manuscript received line
% \IEEELSENSmanuscriptreceived{Manuscript received June 7, 2017;
% revised June 21, 2017; accepted July 6, 2017.
% Date of publication July 12, 2017; date of current version July 12, 2017.}


% for Sensors Letters papers, we must declare the abstract and index terms
% PRIOR to the title within the \IEEEtitleabstractindextext IEEEtran
% command as these need to go into the title area created by \maketitle.
% As a general rule, do not put math, special symbols or citations
% in the abstract or keywords.
% \IEEEtitleabstractindextext{%
% \begin{abstract}
% Direct magnetic field sensors can address integration drift in conventional inductive magnetic diagnostics for fusion. In this work, an AlGaN/GaN Hall-effect sensor was fabricated, packaged, and deployed inside the Helically Symmetric eXperiment (HSX)---the first quasi-helically symmetric stellarator, operating with a 1 T on-axis magnetic field and up to 200 kW of electron cyclotron resonance heating (ECRH)---for in-situ magnetic field monitoring near the plasma edge. The sensor leverages the high-mobility two-dimensional electron gas (2DEG) formed in the AlGaN/GaN heterostructure for sensitive magnetic field measurement, while the wide-bandgap GaN material system provides thermal robustness for harsh-environment operation. The packaged sensor was voltage-biased at 0.4 V, and the output voltage was read out using a differential amplifier chain with a total gain of 200 V/V and a bandwidth of 1 MHz. During 68 consecutive plasma discharge shots, the sensor remained functional and produced clear transient responses associated with plasma ignition and discharge dynamics. Comparisons between biased and unbiased operation, as well as plasma-discharge and coil-only shots, confirmed that the response originated from the biased Hall-effect sensor element. Furthermore, the sensor output exhibited temporal correlation with the plasma stored energy measured by the HSX diamagnetic loop across high-energy, late-breakdown, and failed-breakdown discharges, despite the two diagnostics probing physically distinct quantities. These results demonstrate the feasibility of AlGaN/GaN Hall-effect sensors for in-situ, real-time magnetic diagnostics in fusion-relevant environments.
% \end{abstract}

%150 words version
\IEEEtitleabstractindextext{%
\begin{abstract}[graphical_abstract]
Direct magnetic field sensors can address integration drift commonly observed in conventional inductive magnetic diagnostics used in fusion systems. In this work, an AlGaN/GaN Hall-effect sensor was fabricated, packaged, and deployed inside the Helically Symmetric eXperiment (HSX)---the first quasi-helically symmetric stellarator, operating with a 1 T on-axis magnetic field and up to 200 kW of launched electron cyclotron resonance heating (ECRH) power---for in-situ magnetic field monitoring near the plasma edge. The sensor leverages the high-mobility two-dimensional electron gas (2DEG) formed in the AlGaN/GaN heterostructure for sensitive magnetic field measurement, while the wide-bandgap GaN material system provides thermal robustness for harsh-environment operation. During 68 consecutive plasma discharge shots, the sensor remained functional and produced clear transient responses associated with plasma ignition and discharge dynamics. Comparisons between biased and unbiased operation, as well as plasma-discharge and coil-only shots, confirmed that the response originated from the biased Hall-effect sensor element. Furthermore, the sensor output exhibited temporal correlation with the plasma stored energy measured by the HSX diamagnetic loop across high-energy, late-breakdown, and failed-breakdown discharges.
\end{abstract}

\begin{IEEEkeywords}
Magnetic sensors, Hall-effect sensors, gallium nitride, plasma diagnostics, stellarators.
\end{IEEEkeywords}}


% If you want to put a publisher's ID mark on the page you can do it like
% this:
\IEEEpubid{1949-307X \copyright\ 2023 IEEE. Personal use is permitted, but republication/redistribution requires IEEE permission.\\
See \url{http://www.ieee.org/publications\_standards/publications/rights/index.html} for more information.}
% Remember, if you use this you must call \IEEEpubidadjcol in the second
% column for its text to clear the IEEEpubid mark.


% make the title area
\maketitle


\section{Introduction}
\IEEEPARstart{M}{agnetic} confinement fusion offers the prospect of clean, carbon-free baseload energy, and is regarded as a critical long-term option for future energy supply~\cite{Ongena2016}. Modern fusion experiments increasingly rely on data-driven modeling and machine-learning--based control, placing growing demands on in-situ plasma diagnostics and the associated instrumentation~\cite{Degrave2022, Anirudh2023}. Magnetic field diagnostics are essential for monitoring plasma position, shape, stability, and stored energy in magnetic confinement fusion devices such as tokamaks and stellarators~\cite{ref1}. As fusion experiments advance toward long-pulse and steady-state operation~\cite{ref2}, in-situ diagnostics must provide reliable, real-time measurements of plasma-generated magnetic fields for equilibrium reconstruction and control. Conventional diagnostic systems with inductive sensors, such as Mirnov coils, recover the magnetic field by integrating the induced voltage proportional to its time derivative. During extended operation, this integration becomes problematic as it accumulates low-frequency noise and radiation-induced offsets, leading to drift and reduced accuracy in tracking magnetic fields~\cite{ref3}, \cite{ref4}.

Hall-effect sensors offer an attractive alternative because they enable direct magnetic field measurement without integration, eliminating drift at the source~\cite{ref5}, \cite{ref6}. However, conventional Hall-effect devices based on silicon, gallium arsenide (GaAs), indium arsenide (InAs), and indium antimonide (InSb) platforms cannot be deployed near the plasma edge because the in-vessel environment combines elevated temperature, vacuum operation, chemical reactivity, and radiation exposure. Gallium nitride (GaN) is a wide-bandgap semiconductor with strong thermal tolerance, making it well suited for harsh-environment electronics and in-situ sensing~\cite{ref7}, \cite{ref8}. In AlGaN/GaN heterostructures, a two-dimensional electron gas (2DEG) forms at the interface~\cite{ref9}. Consequently, the high electron mobility of this 2DEG channel enables high sensitivity ideal for this application~\cite{ref10}. In prior thermal characterization experiments, AlGaN/GaN Hall-effect sensors maintained stable current-scaled sensitivity from room temperature up to \SI{576}{\degreeCelsius}~\cite{ref11}, supporting their potential for extreme-environment magnetic field sensing~\cite{ref12}.

\IEEEpubidadjcol In this work, we demonstrate an AlGaN/GaN Hall-effect sensor for in-situ magnetic field monitoring inside the Helically Symmetric eXperiment (HSX), the first stellarator designed and built with a quasi-helically symmetric magnetic field, operated at a \SI{1}{\tesla} on-axis magnetic field~\cite{ref13}, \cite{ref14}. The Hall-effect sensor was fabricated, packaged for ultra-high-vacuum and high-temperature operation, and deployed near the plasma edge. By comparing biased and unbiased operation, plasma-discharge and coil-only shots, and the sensor output with conventional plasma stored-energy diagnostics, we validate that the packaged AlGaN/GaN Hall-effect sensor can detect plasma-dependent magnetic field dynamics in real time, with a 1 MHz readout bandwidth across 68 consecutive shots.
% \newpage
\section{Sensor Fabrication and Experimental Setup}
\subsection{Fabrication}
The AlGaN/GaN wafer, purchased from NTT Advanced Technology Corporation, consists of a \SI{3.7}{\micro\meter} buffer structure, a \SI{300}{\nano\meter} GaN layer, a \SI{1}{\nano\meter} AlN spacer, and a \SI{22}{\nano\meter} Al$_{0.28}$Ga$_{0.72}$N barrier layer. The Hall-effect sensors were fabricated at the Stanford Nanofabrication Facility. The fabrication process began with a mesa etch of the III-nitride layer to define the active regions, followed by deposition and rapid thermal annealing of a Ti~(\SI{25}{\nano\meter})/Al~(\SI{200}{\nano\meter})/Mo~(\SI{40}{\nano\meter})/Au~(\SI{80}{\nano\meter}) stack at \SI{850}{\degreeCelsius} for \SI{35}{\second} to form Ohmic contacts~\cite{ref11}. A \SI{7}{\nano\meter}-thick Al$_2$O$_3$ passivation layer was then deposited; vias were etched to expose the contacts; and a Ti/Au bond-metal layer was deposited for wire bonding. The wafer was singulated into \SI{5}{\milli\meter} $\times$ \SI{5}{\milli\meter} dies for packaging. Fig.~\ref{fig:fab}(a) shows a cross-sectional schematic of the AlGaN/GaN Hall-effect sensor. The fabricated Hall-effect sensors feature a regular octagonal geometry with an inscribed diameter of \SI{200}{\micro\meter}~\cite{ref10}, as illustrated in Fig.~\ref{fig:fab}(b).
\begin{figure}[!t]
\centering
\includegraphics[width=\columnwidth]{./figures/fig1.pdf}
\caption{Fabricated AlGaN/GaN Hall-effect sensor. (a) Cross-sectional schematic of the passivated AlGaN/GaN Hall-effect sensor showing the 2DEG sensing layer, Ohmic contacts, Al$_2$O$_3$ passivation, and bond metal. (b) Optical image of the fabricated regular octagonal Hall-effect plate with a 200 \textmu m inscribed diameter.}
\label{fig:fab}
\end{figure}
\subsection{Packaging}
The individual die was wire-bonded to a ceramic leadless chip carrier (LCC; Spectrum Semiconductor Materials) using aluminum wire. The die and wire bonds were then encapsulated with epoxy (EPO-TEK 353ND) and vacuum-baked at \SI{150}{\degreeCelsius} for 1~hour to meet the ultra-high-vacuum (UHV) requirements of the HSX facility. The packaged sensor was mounted on a custom zirconia ceramic holder attached to a stainless-steel standoff for insertion into the HSX vessel. To reduce the risk of arcing and epoxy degradation during glow discharge cleaning (GDC) and plasma operations, a grounded graphite shield was installed over the packaged sensor module. The packaged sensor assembly and its in-vessel placement are shown in Fig.~\ref{fig:pkg}.
\begin{figure}[!t]
\centering
\includegraphics[width=\columnwidth]{./figures/fig2.pdf}
\caption{Packaging and in-vessel deployment of the AlGaN/GaN Hall-effect sensor in HSX. The sensor die was wire-bonded to a ceramic leadless chip carrier, epoxy-coated, and mounted on a custom flange for insertion into the HSX vessel near the plasma edge.}
\label{fig:pkg}
\end{figure}
\begin{figure}
\centering
\includegraphics[width=\columnwidth]{./figures/fig3.pdf}
\caption{Experimental readout configuration for the in-vessel Hall-effect sensor. The sensor was voltage-biased using the built-in waveform generator of an oscilloscope, and the transverse Hall voltage was amplified using a differential readout chain before data acquisition.}
\label{fig:readout}
\end{figure}
\subsection{Experimental Setup}
The Hall-effect sensor module was deployed inside the HSX vessel near the plasma edge, as shown in Fig.~\ref{fig:pkg}. The bias supply, readout electronics, and data acquisition system were located outside the vessel and connected to the sensor through vacuum feedthroughs. During operation, the sensor was voltage-biased and generated a transverse Hall-effect voltage given by
\begin{equation}
V_{H} = S_{v} V_{\mathrm{bias}} B,
\end{equation}
where $S_{v}$ is the voltage-scaled sensitivity, $B$ is the orthogonal magnetic field, and $V_{\mathrm{bias}}$ is the bias voltage. As illustrated in Fig.~\ref{fig:readout}, two opposite terminals of the Hall-effect plate were biased using the built-in waveform generator of an oscilloscope (Keysight DSOX1204G), while the Hall-effect voltage was measured across the orthogonal terminals. The differential Hall-effect voltage was amplified by a readout circuit consisting of an instrumentation amplifier (Texas Instruments INA849) followed by two operational amplifier (Texas Instruments OPA814) stages, providing a total voltage gain of \SI[per-mode=symbol]{200}{\volt\per\volt} and a bandwidth of \SI{1}{\mega\hertz}. The amplified output voltage can be expressed as
\begin{equation}
V_{\mathrm{out}} = A_{v} V_{H} + V_{\mathrm{off}},
\end{equation}
where $A_{v}$ is the voltage gain and $V_{\mathrm{off}}$ accounts for offset contributions from the sensor and amplifier stages. The amplified output was continuously monitored during the plasma discharge shots. In this readout chain, $V_{\mathrm{off}}$ is bias-independent but may vary with temperature during operation~\cite{ref15}. Since $V_{H}$ scales linearly with $V_{\mathrm{bias}}$ whereas $V_{\mathrm{off}}$ does not, comparing the sensor output under biased ($V_{\mathrm{bias}} = \SI{0.4}{\volt}$) and unbiased ($V_{\mathrm{bias}} = \SI{0}{\volt}$) operation isolates the bias-dependent Hall-effect response from the offset and from other artifacts such as electromagnetic pickup or charge-induced response. Absolute calibration of $V_{\mathrm{off}}$ and quantitative correction for its temperature dependence remain future work.
\begin{figure}[!t]
\centering
\includegraphics[width=\columnwidth]{./figures/fig4.eps}
\caption{Functional verification of the packaged AlGaN/GaN Hall-effect sensor during HSX operation. (a) Amplified sensor output for unbiased (shot 63) and biased (shot 65) operation. (b) Comparison of the biased sensor response during a plasma discharge (shot 65) and a coil-only shot (shot 68).}
\label{fig:func}
\end{figure}
\section{Results and Discussion}
\subsection{Sensor Functionality}
The packaged AlGaN/GaN Hall-effect sensor was deployed inside the HSX stellarator and operated across 68 consecutive shots. Prior to each plasma discharge, mechanical energy stored in 18 motor generators was converted to electrical current driving the 48 magnetic field coils of HSX; after approximately \SI{800}{\milli\second} the coils' current saturated, establishing the target magnetic confinement field. Plasma was then ignited within the resulting \SI{50}{\milli\second} flat-top period using up to \SI{200}{\kilo\watt} of \SI{28}{\giga\hertz} electron cyclotron resonance heating (ECRH). Fig.~\ref{fig:func}(a) compares the amplified output of the Hall-effect sensor when unbiased ($V_{\mathrm{bias}} = \SI{0}{\volt}$; shot 63) and when biased at $V_{\mathrm{bias}} = \SI{0.4}{\volt}$ (shot 65). Consistent with the Hall-effect voltage dependence on $V_{\mathrm{bias}}$, the unbiased output remains near a low baseline, whereas the biased output increases during coil ramp-up and shows a sharp transient followed by fluctuations during the plasma discharge. To separate the plasma-induced response from the background coil field, Fig.~\ref{fig:func}(b) compares a plasma discharge (shot 65) with a coil-only shot (shot 68). Both traces show a similar increase during coil ramp-up, but only the plasma discharge exhibits the pronounced transient and subsequent fluctuations near plasma ignition. These results confirm that the packaged sensor remained functional in the HSX environment and detected magnetic field variations associated with plasma dynamics, with the ignition transient consistent with the diamagnetic response of the confined plasma as plasma pressure rises and expels magnetic flux. Across the 68 discharge shots, the shape and timing of this transient repeated reliably under comparable conditions, indicating stable sensor operation throughout the deployment.
\begin{figure}[!t]
\centering
\includegraphics[width=\columnwidth]{./figures/fig5.eps}
\caption{Real-time tracking of plasma-dependent magnetic-field dynamics. Plasma stored energy from conventional HSX diagnostics is compared with the amplified AlGaN/GaN Hall-effect sensor output for (a), (a$'$)~a high-energy discharge (shot 21); (b), (b$'$)~a late-breakdown discharge (shot 18); and (c), (c$'$)~a failed-breakdown case (shot 19). The two data acquisition systems have a $\sim$30 ms relative timing offset due to delayed triggering.}
\label{fig:track}
\end{figure}
\subsection{Real-Time Plasma Energy Tracking}
The measurements shown in Fig.~\ref{fig:track} compare the amplified output of the in-vessel AlGaN/GaN Hall-effect sensor with the plasma stored energy measured by the HSX diamagnetic loop diagnostic~\cite{ref16} for three representative discharge conditions. The two diagnostic methods measure physically distinct quantities: the Hall-effect sensor probes the local magnetic field at a single in-vessel location, whereas the diamagnetic loop measures the change in enclosed toroidal flux from which volumetric plasma stored energy is inferred. Both, however, respond to plasma magnetic dynamics and are therefore expected to be temporally correlated~\cite{ref1}. The two signals are acquired through separate data acquisition systems with a known $\sim$\SI{30}{\milli\second} relative timing offset arising from delayed triggering of the diamagnetic loop DAQ. In the high-energy discharge (shot 21) shown in Fig.~\ref{fig:track}(a) and (a$'$), the rapid rise and sustained stored-energy plateau are accompanied by a pronounced sensor response and elevated fluctuations during the discharge. In the late-breakdown case (shot 18) shown in Fig.~\ref{fig:track}(b) and (b$'$), the sensor output remains near baseline until the delayed rise in plasma stored energy, after which a clear increase in signal amplitude is observed. In contrast, for the failed-breakdown case (shot 19) shown in Fig.~\ref{fig:track}(c) and (c$'$), both the stored energy and the sensor output remain near baseline. Although quantitative calibration remains future work, the consistent temporal agreement across these discharge scenarios establishes that the AlGaN/GaN Hall-effect sensor can track plasma-dependent magnetic-field dynamics in real time.
\section{Conclusion}
An AlGaN/GaN Hall-effect sensor was fabricated, packaged, and deployed inside the HSX stellarator (\SI{1}{\tesla} on-axis field, up to \SI{200}{\kilo\watt} of ECRH) for in-situ magnetic-field monitoring near the plasma edge. The packaged device operated through 68 consecutive shots under ultra-high-vacuum and elevated-temperature conditions. Comparisons between biased and unbiased operation, together with plasma-discharge and coil-only measurements, confirmed that the observed response was associated with magnetic field changes detected through the biased Hall-effect sensor. The amplified sensor output further showed temporal agreement with conventional plasma stored-energy diagnostics across high-energy, late-breakdown, and failed-breakdown scenarios. Although quantitative calibration remains future work, these results establish AlGaN/GaN Hall-effect sensors as a feasible platform for real-time magnetic diagnostics in fusion-relevant environments. Future work will focus on absolute calibration through in-situ cross-reference with established HSX magnetic diagnostics~\cite{ref17}, extended-duration deployments in HSX to evaluate offset stability, radiation and neutron irradiation characterization at a dedicated facility to assess radiation hardness, and integration of lower-noise readout electronics for resolving smaller-amplitude magnetohydrodynamic fluctuations.

% \section*{APPENDIX}

% Appendixes, if needed, appear before the acknowledgment.
\section*{Acknowledgment}
\addcontentsline{toc}{section}{Acknowledgment}
\scriptsize
The work of the authors was supported by the U.S.\ Department of Energy under Contract No.~DE-AC02-76SF00515, SLAC FWP~101264, and by the TomKat Center for Sustainable Energy at Stanford University. Fabrication work was performed at the Stanford Nanofabrication Facility, a member of the National Nanotechnology Coordinated Infrastructure (NNCI), supported by the National Science Foundation under Award~ECCS-2026822.
% put at least one blank line to end the scriptsize paragraph and
% then revert back to normalsize.
\normalsize


% Last page column equalization
%
% IEEE Sensors Letters does balance the columns on the last page.
% Can use:
% \IEEEtriggeratref{8}
% to trigger a \newpage just before the given reference number to
% balance the columns on the last page. Adjust the reference number
% as needed - this may need to be readjusted if the document is 
% modified later.
% The "triggered" command can be changed if desired:
%\IEEEtriggercmd{\enlargethispage{-5in}}
%
% Alternatively, you can also directly use something like
% \enlargethispage{-7in}
% on the last page instead of breaking at a specific reference number.


% references section
%
% can use a bibliography generated by BibTeX as a .bbl file
% BibTeX documentation can be easily obtained at:
% http://mirror.ctan.org/biblio/bibtex/contrib/doc/
% The IEEEtran BibTeX style support page is at:
% http://www.michaelshell.org/tex/ieeetran/bibtex/
%\bibliographystyle{IEEEtran}
% argument is your BibTeX string definitions and bibliography database(s)
%\bibliography{IEEEabrv,../bib/paper}
%
% Before submitting to IEEE Sensors Letters, manually copy in the
% resultant .bbl file contents in place of the \bibliographystyle and
% \bibliography lines here:
\begin{thebibliography}{20}

\bibitem{Ongena2016}
J. Ongena, R. Koch, R. Wolf, and H. Zohm, ``Magnetic-confinement fusion,'' \emph{Nat. Phys.}, vol. 12, no. 5, pp. 398--410, May 2016, doi: 10.1038/nphys3745.

\bibitem{Degrave2022}
J. Degrave \emph{et al.}, ``Magnetic control of tokamak plasmas through deep reinforcement learning,'' \emph{Nature}, vol. 602, no. 7897, pp. 414--419, Feb. 2022, doi: 10.1038/s41586-021-04301-9.

\bibitem{Anirudh2023}
R. Anirudh \emph{et al.}, ``2022 review of data-driven plasma science,'' \emph{IEEE Trans. Plasma Sci.}, vol. 51, no. 7, pp. 1750--1838, Jul. 2023, doi: 10.1109/TPS.2023.3268170.

\bibitem{ref1}
E. J. Strait, E. D. Fredrickson, J.-M. Moret, and M. Takechi, ``Chapter 2: Magnetic diagnostics,'' \emph{Fusion Sci. Technol.}, vol. 53, no. 2, pp. 304--334, Feb. 2008, doi: 10.13182/FST08-A1674.

\bibitem{ref2}
M. Endler \emph{et al.}, ``Engineering design for the magnetic diagnostics of Wendelstein 7-X,'' \emph{Fusion Eng. Des.}, vol. 100, pp. 468--494, Nov. 2015, doi: 10.1016/j.fusengdes.2015.07.020.

\bibitem{ref3}
I. \v{D}uran \emph{et al.}, ``Status of steady-state magnetic diagnostic for ITER and outlook for possible materials of Hall sensors for DEMO,'' \emph{Fusion Eng. Des.}, vol. 146, pp. 2397--2400, Sep. 2019, doi: 10.1016/j.fusengdes.2019.03.201.

\bibitem{ref4}
E. J. Strait, ``Magnetic diagnostic system of the DIII-D tokamak,'' \emph{Rev. Sci. Instrum.}, vol. 77, no. 2, Art. no. 023502, Feb. 2006, doi: 10.1063/1.2166493.

\bibitem{ref5}
A. Quercia \emph{et al.}, ``Long term operation of the radiation-hard Hall probes system and the path toward a high performance hybrid magnetic field sensor,'' \emph{Nucl. Fusion}, vol. 62, no. 10, Art. no. 106032, Sep. 2022, doi: 10.1088/1741-4326/ac8aad.

\bibitem{ref6}
I. Bolshakova \emph{et al.}, ``Metal Hall sensors for the new generation fusion reactors of DEMO scale,'' \emph{Nucl. Fusion}, vol. 57, no. 11, Art. no. 116042, Aug. 2017, doi: 10.1088/1741-4326/aa7867.

\bibitem{ref7}
U. K. Mishra, P. Parikh, and Y.-F. Wu, ``AlGaN/GaN HEMTs---an overview of device operation and applications,'' \emph{Proc. IEEE}, vol. 90, no. 6, pp. 1022--1031, Jun. 2002, doi: 10.1109/JPROC.2002.1021567.

\bibitem{ref8}
M. Yuan \emph{et al.}, ``Enhancement-mode GaN transistor technology for harsh environment operation,'' \emph{IEEE Electron Device Lett.}, vol. 44, no. 7, pp. 1068--1071, Jul. 2023, doi: 10.1109/LED.2023.3279813.

\bibitem{ref9}
O. Ambacher \emph{et al.}, ``Two-dimensional electron gases induced by spontaneous and piezoelectric polarization charges in N- and Ga-face AlGaN/GaN heterostructures,'' \emph{J. Appl. Phys.}, vol. 85, no. 6, pp. 3222--3233, Mar. 1999, doi: 10.1063/1.369664.

\bibitem{ref10}
H. S. Alpert \emph{et al.}, ``Effect of geometry on sensitivity and offset of AlGaN/GaN and InAlN/GaN Hall-effect sensors,'' \emph{IEEE Sensors J.}, vol. 19, no. 10, pp. 3640--3646, May 2019, doi: 10.1109/JSEN.2019.2895546.

\bibitem{ref11}
H. S. Alpert \emph{et al.}, ``Sensitivity of 2DEG-based Hall-effect sensors at high temperatures,'' \emph{Rev. Sci. Instrum.}, vol. 91, no. 2, Art. no. 025003, Feb. 2020, doi: 10.1063/1.5139911.

\bibitem{ref12}
H. Lu, P. Sandvik, A. Vertiatchikh, J. Tucker, and A. Elasser, ``High temperature Hall effect sensors based on AlGaN/GaN heterojunctions,'' \emph{J. Appl. Phys.}, vol. 99, no. 11, Art. no. 114510, Jun. 2006, doi: 10.1063/1.2201339.

\bibitem{ref13}
F. S. B. Anderson, A. F. Almagri, D. T. Anderson, P. G. Matthews, J. N. Talmadge, and J. L. Shohet, ``The Helically Symmetric Experiment, (HSX) goals, design and status,'' \emph{Fusion Technol.}, vol. 27, no. 3T, pp. 273--277, Apr. 1995, doi: 10.13182/FST95-A11947086.

\bibitem{ref14}
A. F. Almagri \emph{et al.}, ``Design and construction of HSX: A helically symmetric stellarator,'' \emph{J. Plasma Fusion Res. SERIES}, vol. 1, pp. 422--425, 1998.

\bibitem{ref15}
K. M. Dowling, ``Offset and noise behavior of microfabricated aluminum gallium nitride--gallium nitride two-dimensional electron gas Hall-effect sensors,'' Ph.D. dissertation, Dept. Elect. Eng., Stanford Univ., Stanford, CA, USA, 2019.

\bibitem{ref16}
J. C. Schmitt, J. N. Talmadge, and D. T. Anderson, ``Measurement of a helical Pfirsch--Schl\"uter current with reduced magnitude in HSX,'' \emph{Nucl. Fusion}, vol. 53, no. 8, Art. no. 082001, Jul. 2013, doi: 10.1088/0029-5515/53/8/082001.

\bibitem{ref17}
E. Chlechowitz, J. N. Talmadge, J. D. Hanson, F. S. B. Anderson, and D. T. Anderson, ``Magnetic diagnostic optimization for plasma equilibrium reconstruction in the HSX stellarator,'' \emph{Nucl. Fusion}, vol. 55, no. 11, Art. no. 113012, Sep. 2015, doi: 10.1088/0029-5515/55/11/113012.

\end{thebibliography}


% that's all folks
\end{document}


